% Options for packages loaded elsewhere
\PassOptionsToPackage{unicode}{hyperref}
\PassOptionsToPackage{hyphens}{url}
%
\documentclass[
]{article}
\usepackage{amsmath,amssymb}
\usepackage{lmodern}
\usepackage{iftex}
\ifPDFTeX
  \usepackage[T1]{fontenc}
  \usepackage[utf8]{inputenc}
  \usepackage{textcomp} % provide euro and other symbols
\else % if luatex or xetex
  \usepackage{unicode-math}
  \defaultfontfeatures{Scale=MatchLowercase}
  \defaultfontfeatures[\rmfamily]{Ligatures=TeX,Scale=1}
\fi
% Use upquote if available, for straight quotes in verbatim environments
\IfFileExists{upquote.sty}{\usepackage{upquote}}{}
\IfFileExists{microtype.sty}{% use microtype if available
  \usepackage[]{microtype}
  \UseMicrotypeSet[protrusion]{basicmath} % disable protrusion for tt fonts
}{}
\makeatletter
\@ifundefined{KOMAClassName}{% if non-KOMA class
  \IfFileExists{parskip.sty}{%
    \usepackage{parskip}
  }{% else
    \setlength{\parindent}{0pt}
    \setlength{\parskip}{6pt plus 2pt minus 1pt}}
}{% if KOMA class
  \KOMAoptions{parskip=half}}
\makeatother
\usepackage{xcolor}
\usepackage[margin=1in]{geometry}
\usepackage{longtable,booktabs,array}
\usepackage{calc} % for calculating minipage widths
% Correct order of tables after \paragraph or \subparagraph
\usepackage{etoolbox}
\makeatletter
\patchcmd\longtable{\par}{\if@noskipsec\mbox{}\fi\par}{}{}
\makeatother
% Allow footnotes in longtable head/foot
\IfFileExists{footnotehyper.sty}{\usepackage{footnotehyper}}{\usepackage{footnote}}
\makesavenoteenv{longtable}
\usepackage{graphicx}
\makeatletter
\def\maxwidth{\ifdim\Gin@nat@width>\linewidth\linewidth\else\Gin@nat@width\fi}
\def\maxheight{\ifdim\Gin@nat@height>\textheight\textheight\else\Gin@nat@height\fi}
\makeatother
% Scale images if necessary, so that they will not overflow the page
% margins by default, and it is still possible to overwrite the defaults
% using explicit options in \includegraphics[width, height, ...]{}
\setkeys{Gin}{width=\maxwidth,height=\maxheight,keepaspectratio}
% Set default figure placement to htbp
\makeatletter
\def\fps@figure{htbp}
\makeatother
\setlength{\emergencystretch}{3em} % prevent overfull lines
\providecommand{\tightlist}{%
  \setlength{\itemsep}{0pt}\setlength{\parskip}{0pt}}
\setcounter{secnumdepth}{-\maxdimen} % remove section numbering
\newlength{\cslhangindent}
\setlength{\cslhangindent}{1.5em}
\newlength{\csllabelwidth}
\setlength{\csllabelwidth}{3em}
\newlength{\cslentryspacingunit} % times entry-spacing
\setlength{\cslentryspacingunit}{\parskip}
\newenvironment{CSLReferences}[2] % #1 hanging-ident, #2 entry spacing
 {% don't indent paragraphs
  \setlength{\parindent}{0pt}
  % turn on hanging indent if param 1 is 1
  \ifodd #1
  \let\oldpar\par
  \def\par{\hangindent=\cslhangindent\oldpar}
  \fi
  % set entry spacing
  \setlength{\parskip}{#2\cslentryspacingunit}
 }%
 {}
\usepackage{calc}
\newcommand{\CSLBlock}[1]{#1\hfill\break}
\newcommand{\CSLLeftMargin}[1]{\parbox[t]{\csllabelwidth}{#1}}
\newcommand{\CSLRightInline}[1]{\parbox[t]{\linewidth - \csllabelwidth}{#1}\break}
\newcommand{\CSLIndent}[1]{\hspace{\cslhangindent}#1}
\ifLuaTeX
  \usepackage{selnolig}  % disable illegal ligatures
\fi
\IfFileExists{bookmark.sty}{\usepackage{bookmark}}{\usepackage{hyperref}}
\IfFileExists{xurl.sty}{\usepackage{xurl}}{} % add URL line breaks if available
\urlstyle{same} % disable monospaced font for URLs
\hypersetup{
  pdftitle={Executive Summary of Cook Inlet Beluga eDNA Results},
  pdfauthor={Zack Gold},
  hidelinks,
  pdfcreator={LaTeX via pandoc}}

\title{Executive Summary of Cook Inlet Beluga eDNA Results}
\author{Zack Gold}
\date{2023-11-27}

\begin{document}
\maketitle

\textbf{Note: These are preliminary data and a first look at the eDNA
results. Please only share internally with other NOAA \& AKBMP partners
and collaborators.}

\hypertarget{project-overview}{%
\section{Project Overview}\label{project-overview}}

These are the results of a pilot study of eDNA in Cook Inlet from
Nov.~2020 to Nov.~2021. This pilot project consisted of 3 eDNA sample
collections:

\begin{enumerate}
\def\labelenumi{\arabic{enumi}.}
\item
  Monthly samples were collected at two sites in the Kenai River across
  the year,
\item
  Higher frequency sampling (weekly/biweelky) of 20 Mile River was
  conducted from August 2021-Nov.~2021, and
\item
  Oppurtunistic sampling within and adjacent to Belugas from ADF\&G
  small vessels in August 2021.
\end{enumerate}

We successfully designed, optimized, and implemented sample collection
and processing methodologies to track Beluga and their prey. We
specifically applied qPCR assays for Beluga, Coho, Chinook, and Eulachon
(\emph{Thaleichthys pacificus}) and vertebrate community metabarcoding
assays (MiFish 12S) to all eDNA samples.

\hypertarget{summary-of-major-findings-results}{%
\section{Summary of Major Findings
Results}\label{summary-of-major-findings-results}}

\hypertarget{mile-river}{%
\section{20 Mile River}\label{mile-river}}

\hypertarget{edna-metabarcoding-vs.-qpcr-vs.-visual-observations}{%
\subsection{eDNA Metabarcoding vs.~qPCR vs.~Visual
Observations}\label{edna-metabarcoding-vs.-qpcr-vs.-visual-observations}}

eDNA and visual observations of Cook Inlet Beluga whales were largely
concordant within Turnagain Arm (Figure 1). eDNA detections at 20 Mile
River during periods of time in which Beluga were only visually observed
at Bird Point strongly suggest that Beluga DNA is being transported
\textgreater10 km. These results demonstrate that both metabarcoding and
qPCR eDNA approaches can provide valuable data on the presence of
beluga.

\hypertarget{figure-1}{%
\subsubsection{Figure 1}\label{figure-1}}

\includegraphics{CIB_analysis_exec_sum_20231126_files/figure-latex/unnamed-chunk-2-1.pdf}

\hypertarget{beluga-run-timing-strongly-correlated-with-coho-salmon}{%
\subsection{Beluga run timing strongly correlated with Coho
salmon}\label{beluga-run-timing-strongly-correlated-with-coho-salmon}}

We demonstrate that eDNA metabarcoding can successfully capture changes
in fish assemblages across both space and time allowing for comparison
to Beluga foraging patterns. Specifically, we find that Beluga run
timing was strongly correlated with increases in Coho abundances (Figure
2).

\hypertarget{figure-2}{%
\section{Figure 2}\label{figure-2}}

We see similar patterns when comparing across eDNA indices (Figure 3).

\#Figure 3

\includegraphics{CIB_analysis_exec_sum_20231126_files/figure-latex/unnamed-chunk-3-1.pdf}

From the most abundantly detected taxa (over 1,000 reads detected in at
least 10 technical replicates), we found signficant differences in eDNA
index scores before and during beluga occurrences in 20 Mile River.

\hypertarget{more-abundant-prior-to-beluga-occurrence}{%
\paragraph{More Abundant Prior to Beluga
Occurrence}\label{more-abundant-prior-to-beluga-occurrence}}

\begin{longtable}[]{@{}
  >{\raggedright\arraybackslash}p{(\columnwidth - 10\tabcolsep) * \real{0.3125}}
  >{\raggedright\arraybackslash}p{(\columnwidth - 10\tabcolsep) * \real{0.0875}}
  >{\raggedleft\arraybackslash}p{(\columnwidth - 10\tabcolsep) * \real{0.2000}}
  >{\raggedleft\arraybackslash}p{(\columnwidth - 10\tabcolsep) * \real{0.1750}}
  >{\raggedleft\arraybackslash}p{(\columnwidth - 10\tabcolsep) * \real{0.1000}}
  >{\raggedleft\arraybackslash}p{(\columnwidth - 10\tabcolsep) * \real{0.1250}}@{}}
\toprule()
\begin{minipage}[b]{\linewidth}\raggedright
Species
\end{minipage} & \begin{minipage}[b]{\linewidth}\raggedright
Time
\end{minipage} & \begin{minipage}[b]{\linewidth}\raggedleft
Avg. eDNA Index
\end{minipage} & \begin{minipage}[b]{\linewidth}\raggedleft
SD eDNA Index
\end{minipage} & \begin{minipage}[b]{\linewidth}\raggedleft
p.value
\end{minipage} & \begin{minipage}[b]{\linewidth}\raggedleft
statistic
\end{minipage} \\
\midrule()
\endhead
Oncorhynchus gorbuscha & Before & 0.6285946 & 0.1467596 & 0 &
60.30885 \\
Oncorhynchus gorbuscha & During & 0.2973284 & 0.2942747 & 0 &
60.30885 \\
Oncorhynchus tshawytscha & Before & 0.1797197 & 0.1253640 & 0 &
72.23416 \\
Oncorhynchus tshawytscha & During & 0.0140808 & 0.0169829 & 0 &
72.23416 \\
Thaleichthys pacificus & Before & 0.1331993 & 0.1150512 & 0 &
37.74665 \\
Thaleichthys pacificus & During & 0.0217129 & 0.0294135 & 0 &
37.74665 \\
\bottomrule()
\end{longtable}

\hypertarget{more-abundant-during-beluga-occurrence}{%
\paragraph{More Abundant During Beluga
Occurrence}\label{more-abundant-during-beluga-occurrence}}

\begin{longtable}[]{@{}
  >{\raggedright\arraybackslash}p{(\columnwidth - 10\tabcolsep) * \real{0.2692}}
  >{\raggedright\arraybackslash}p{(\columnwidth - 10\tabcolsep) * \real{0.0897}}
  >{\raggedleft\arraybackslash}p{(\columnwidth - 10\tabcolsep) * \real{0.2051}}
  >{\raggedleft\arraybackslash}p{(\columnwidth - 10\tabcolsep) * \real{0.1795}}
  >{\raggedleft\arraybackslash}p{(\columnwidth - 10\tabcolsep) * \real{0.1282}}
  >{\raggedleft\arraybackslash}p{(\columnwidth - 10\tabcolsep) * \real{0.1282}}@{}}
\toprule()
\begin{minipage}[b]{\linewidth}\raggedright
Species
\end{minipage} & \begin{minipage}[b]{\linewidth}\raggedright
Time
\end{minipage} & \begin{minipage}[b]{\linewidth}\raggedleft
Avg. eDNA Index
\end{minipage} & \begin{minipage}[b]{\linewidth}\raggedleft
SD eDNA Index
\end{minipage} & \begin{minipage}[b]{\linewidth}\raggedleft
p.value
\end{minipage} & \begin{minipage}[b]{\linewidth}\raggedleft
statistic
\end{minipage} \\
\midrule()
\endhead
Oncorhynchus kisutch & Before & 0.0585469 & 0.0372569 & 0.0000000 &
64.169594 \\
Oncorhynchus kisutch & During & 0.4243747 & 0.3690432 & 0.0000000 &
64.169594 \\
Oncorhynchus nerka & Before & 0.0213065 & 0.0115372 & 0.0000000 &
39.133972 \\
Oncorhynchus nerka & During & 0.0946590 & 0.0944066 & 0.0000000 &
39.133972 \\
Pungitius pungitius & Before & 0.0405959 & 0.0513276 & 0.0028031 &
9.364055 \\
Pungitius pungitius & During & 0.0752271 & 0.0657437 & 0.0028031 &
9.364055 \\
Salvelinus malma & Before & 0.0299390 & 0.0292162 & 0.0067172 &
7.644705 \\
Salvelinus malma & During & 0.0475393 & 0.0365466 & 0.0067172 &
7.644705 \\
\bottomrule()
\end{longtable}

Coho is a known prey item for Cook Inlet Belugas (Quakenbush et al.
2015) and thus aligns well with previous research (McGuire et al. 2020;
Nelson et al. 2018; Castellote et al. 2020). Traditional ecological
knowledge from informal interviews with local fisherman additionally
supported that the timing of the Coho run coincided with the timing of
Belugas arriving in 20 Mile river.

\hypertarget{kenai-river}{%
\section{Kenai River}\label{kenai-river}}

\hypertarget{visual-observations}{%
\subsection{Visual Observations}\label{visual-observations}}

\hypertarget{figure-4}{%
\subsubsection{Figure 4}\label{figure-4}}

\includegraphics{CIB_analysis_exec_sum_20231126_files/figure-latex/unnamed-chunk-7-1.pdf}

Visual observations of Belugas in Kenai River demonstrate a noted
absence in the summer with presence in Fall and Spring during AKBMP
observations (Figure 4). Here we focus on what fish species were
observed in higher abundances when Beluga were present by isolating the
Winter to Spring transition and the Summer to Fall transition. In both
cases Beluga are conspicously absent from the river in Winter/Summer,
but present in Spring (March 25-April 30th, 2021) and Fall (September
1st-November 2021).

\hypertarget{winter-to-spring}{%
\subsection{Winter to Spring}\label{winter-to-spring}}

\hypertarget{summer-to-fall}{%
\subsection{Summer to Fall}\label{summer-to-fall}}

\hypertarget{conclusions}{%
\subsection{Conclusions}\label{conclusions}}

We demonstrate that eDNA metabarcoding can successfully capture changes
in fish assemblages across both space and time allowing for comparison
to Beluga foraging patterns. Time of sampling across seasons explained
the greatest proportion of variance in our fish assemblage data aligning
well with expectations of seasonally influenced fish assemblages driven
by timing of anadromous migrations. We found that eDNA approaches are
able to capture changes in fish assemblages on the sub-weekly scale,
allowing for high resolution tracking of run timing. Site also had a
signficant effect on fish assemblages, particularly across Cook Inlet
rivers with less pronounced, but signficant differences across sites
along the Kenai River. We observed no statistically signficant
difference in depth within the Kenai River. Together, these results
provide important baseline information needed to inform effective eDNA
sampling efforts needed to resolve Cook Inlet Beluga foraging patterns.
Specifically, these findings suggest that weekly sampling efforts at a
single river site along multiple important foraging habitats would
provide sufficient resolution to identify spatial and temporal river
foraging habitat utilization of Cook Inlet Beluga whales.

Ultimately, these results demonstrate that eDNA approaches are an
effective tool for capturing changes of Beluga prey assemblages across
both space and time. Targeted deployments of eDNA across Cook Inlet
habits can reveal critical ecological insights into fish ecology and
endangered Cook Inlet Beluga foraging patterns.

\hypertarget{references}{%
\section*{References}\label{references}}
\addcontentsline{toc}{section}{References}

\hypertarget{refs}{}
\begin{CSLReferences}{1}{0}
\leavevmode\vadjust pre{\hypertarget{ref-castellote_seasonal_2020}{}}%
Castellote, Manuel, Robert J. Small, Marc O. Lammers, Justin Jenniges,
Jeffrey Mondragon, Christopher D. Garner, Shannon Atkinson, Jade MS
Delevaux, Richard Graham, and Delmar Westerholt. 2020. {``Seasonal
Distribution and Foraging Occurrence of Cook Inlet Beluga Whales Based
on Passive Acoustic Monitoring.''} \emph{Endangered Species Research}
41: 225--43. \url{https://www.int-res.com/abstracts/esr/v41/p225-243}.

\leavevmode\vadjust pre{\hypertarget{ref-mcguire_distribution_2020}{}}%
McGuire, Tamara L., Gina K. Himes Boor, John R. McClung, Amber D.
Stephens, Christopher Garner, Kim E. W. Shelden, and Bruce Wright. 2020.
{``Distribution and Habitat Use by Endangered Cook Inlet Beluga Whales:
Patterns Observed During a Photo‐identification Study, 2005--2017.''}
\emph{Aquatic Conservation: Marine and Freshwater Ecosystems} 30 (12):
2402--27. \url{https://doi.org/10.1002/aqc.3378}.

\leavevmode\vadjust pre{\hypertarget{ref-nelson_fifty_2018}{}}%
Nelson, Mark A., Lori T. Quakenbush, Barbara A. Mahoney, Brian D. Taras,
and Matthew J. Wooller. 2018. {``Fifty Years of Cook Inlet Beluga Whale
Feeding Ecology from Isotopes in Bone and Teeth.''} \emph{Endangered
Species Research} 36: 77--87.
\url{https://www.int-res.com/abstracts/esr/v36/p77-87/}.

\leavevmode\vadjust pre{\hypertarget{ref-quakenbush_diet_2015}{}}%
Quakenbush, Lori T., Robert S. Suydam, Anna L. Bryan, Lloyd F. Lowry,
Kathryn J. Frost, and Barbara A. Mahoney. 2015. {``Diet of Beluga Whales
(Delphinapterus Leucas) in Alaska from Stomach Contents,
March--November.''} \emph{Mar Fish Rev} 77 (1): 70--84.
\url{https://spo.nmfs.noaa.gov/sites/default/files/pdf-content/MFR/mfr771/mfr7717.pdf}.

\end{CSLReferences}

\end{document}
